\section{Experiments}
\label{sec:experiments}
\begin{table*}[t]
\centering\small
\caption{Efficiency of dense and PairSelect configurations. $F_{\rm blocks}$ is an analytical principal-block estimate (1 MAC = 2 FLOPs). Full-encoder FLOPs, latency, throughput, and peak allocated/reserved memory remain unmeasured.}
\label{tab:efficiency}
\setlength{\tabcolsep}{5.5pt}
\begin{tabular}{@{}llrrrrrr@{}}
\toprule
Encoder & Delete & Tokens & $F_{\rm blocks}$ & $F_{\rm full}$ & Latency & Throughput & Peak memory\\
 & & & GFLOPs & GFLOPs & ms/clip & clips/s & alloc./res. GiB\\
\midrule
\tableinput{tables/efficiency_rows}
\bottomrule
\end{tabular}
\end{table*}

\subsection{Four-system evaluation}
We evaluate on UCF-Crime~\cite{sultani2018real} (1,610 training and 290 test videos) and XD-Violence~\cite{wu2020xd} (3,950 accepted training and all 800 test videos; four declared training files fail source-integrity checks). VideoMAEv2-B, VideoMAE-B, and TimeSformer-B use pretrained encoders and separate UR-DMU heads trained on the corresponding complete accepted training sets. The heads use 200 temporal bins, 64 normal and 64 anomalous bags per step, Adam with learning rate $10^{-4}$ and weight decay $5\!\times\!10^{-5}$, and the final checkpoint after 3,000 steps. CLIP ViT-B/16 uses the authors' released UCF/XD DSANet checkpoints~\cite{yin2026learning}, with no DSANet training in this study. All four systems reuse their dense detector at every budget.

Each system preserves its native input schedule and temporal scoring protocol. VideoMAEv2-B and VideoMAE-B use 16-frame inputs, TimeSformer-B uses eight-frame inputs, and CLIP ViT-B/16 processes sampled frames under its fixed raw-video protocol, all at 224$\times$224 resolution. Detector weights and scoring rules remain fixed across token budgets, with the coarse branch used for DSANet. The CLIP dense reference is extracted from the same raw videos as its compressed counterparts, ensuring that the comparison uses the evaluated raw-video pipeline rather than a score obtained from externally released features.

UCF uses frame ROC-AUC as primary metric and also reports step AP. On XD, the three video systems use trapezoidal PR-AUC with the fixed 16-frame prefix alignment, while DSANet uses step AP. Table~\ref{tab:quality} separates both integration definitions. Paired video-level, weak-label-stratified bootstrap uses 10,000 resamples; all-budget intervals are retained with the result exports.

\subsection{Quality across systems and budgets}
At 40\% nominal deletion, \method{} is nearly lossless in most configurations. The three video systems change UCF AUC by $+0.39$, $-0.16$, and approximately $0.00$ points, and XD PR-AUC by $+0.32$, $+0.61$, and $+0.68$. DSANet changes UCF AUC by $-0.90$ and XD step AP by $-0.24$ points, retaining 87.97 AUC and 85.35 AP. UCF step AP is also stable: 19.29$\to$19.60, 18.82$\to$18.99, 20.37$\to$20.90, and 34.74$\to$34.36 across the four systems. No compression-specific training is used to obtain these results.

Increasing deletion to 60\% generally incurs mild degradation while preserving high detection quality. Several configurations improve, and the three largest primary-metric drops are DSANet/UCF ($-1.80$), DSANet/XD ($-1.22$), and VideoMAE/UCF ($-1.05$ points). This extends the usable reduction range across video-pretrained and image--text-pretrained representations and across UR-DMU and DSANet. The two XD metrics need not move together: VideoMAE at 60\% gains 0.78 PR-AUC points but loses 1.39 step-AP points; both are retained in the table. Near-lossless describes point estimates rather than a blanket statistical guarantee: among DSANet budgets, only XD at 20\% has its 95\% lower bound above the predeclared $-0.5$-point margin.

\subsection{Encoder efficiency}
Table~\ref{tab:efficiency} compares batch-32 encoder efficiency across budgets using XD training inputs. Encoder timing includes embedding, selection, gathering, all blocks, and readout. Video-encoder B32 tests use two fixed training videos, 20 warmups and 40 synchronized samples; CLIP uses three processes, 20 warmups and 40 timed samples per setting. Arithmetic values cover the principal twelve blocks (one MAC is two FLOPs), not a full-model profiler estimate.

We evaluate encoder efficiency using speedup relative to each system’s dense counterpart at the same batch size. At batch size 32, PairSelect achieves encoder speedups of 1.22--1.28$\times$ at 40\% nominal token deletion while preserving near-dense detection performance in most configurations. Increasing deletion to 60\% raises encoder speedups to 1.33--1.47$\times$, with generally mild performance degradation. These results highlight a favorable accuracy--efficiency trade-off: PairSelect reduces internal token computation while preserving input coverage and reusing the same detector weights across budgets.
